% ---------------------------------------------------------------------
%  CHƯƠNG 4 — CÀI ĐẶT VÀ THỬ NGHIỆM
% ---------------------------------------------------------------------
\chapter{Cài đặt và thử nghiệm}
\label{ch:caidat}

\begin{muctieu}
Chương này mô tả môi trường triển khai, trách nhiệm của các thành phần, giao
diện đã cài đặt và kết quả kiểm thử hệ thống.
\end{muctieu}

\section{Môi trường triển khai}

Ứng dụng \code{app\_nhan\_dien} dùng Flutter 3.47.1 stable và Dart 3.13.1.
Mã nguồn có \code{FirebaseOptions} cho Android, iOS, Web và Windows; phạm vi
cài đặt, chạy thử và kiểm thử của báo cáo tập trung vào Android. Firebase
Authentication hỗ trợ xác thực email/mật khẩu~\cite{FirebaseAuthFlutter}.
Cloud Firestore tổ chức dữ liệu theo tài liệu và collection~\cite{FirestoreDataModel};
trong mã dự án, dịch vụ này lưu tài liệu \code{users/\{uid\}} và trường hồ sơ cơ thể. Các tuỳ chọn phân tích và lịch
sử hỏi đáp được lưu cục bộ bằng \code{shared\_preferences} 2.5.5. Thành phần
trí tuệ nhân tạo được truy cập qua FastAPI 0.141.1 chạy bằng Uvicorn 0.52.4.

Thông tin trên Firebase Console được đối chiếu với cấu hình và lời gọi dịch
vụ trong mã nguồn, kết quả trình bày ở Bảng~\ref{tab:cauhinh-firebase}. Cách
đối chiếu này tách biệt dịch vụ đã đăng ký trên Console với dịch vụ thật sự
được ứng dụng gọi trong phiên bản hiện tại.

\begin{table}[H]
\centering
\caption{Cấu hình dự án Firebase tại thời điểm đối chiếu ngày 11/11/2026.}
\label{tab:cauhinh-firebase}
\footnotesize
\begin{tabularx}{\textwidth}{>{\bfseries\color{daunavy}\raggedright\arraybackslash}p{3.4cm} >{\raggedright\arraybackslash}X}
\toprule
\rowcolor{gray100}\bfseries\color{daunavy}Hạng mục & \bfseries\color{daumaroon}Cấu hình và trạng thái sử dụng \\
\midrule
Dự án và gói dịch vụ & Tên hiển thị \code{ai-fashion-project}, Project ID \code{ai-fashion-project-6fa98}; gói Spark (no-cost). \\
Ứng dụng đã đăng ký & Android \code{com.example.app\_nhan\_dien}, iOS \code{com.example.appNhanDien} và hai cấu hình Web. Các định danh này khớp với \code{firebase\_options.dart} và \code{google-services.json}. \\
Authentication & Email/Password đã bật. Mã Flutter thực hiện đăng ký, đăng nhập, đăng xuất, đặt lại và đổi mật khẩu bằng Firebase Auth. Tên miền được uỷ quyền gồm \code{localhost}, \code{ai-fashion-project-6fa98.firebaseapp.com} và \code{ai-fashion-project-6fa98.web.app}. \\
Cloud Firestore & Cơ sở dữ liệu Native đặt tại \code{nam5}. Ứng dụng đọc/ghi hồ sơ tại \code{users/\{uid\}}; Rules chỉ cho người đã xác thực thao tác tài liệu của chính mình và subcollection \code{wardrobe}. \\
Google Analytics & Firebase Console đã liên kết Property ID \code{553061587}; cấu hình Web có Measurement ID. Tuy nhiên, mã Flutter không khai báo \code{firebase\_analytics} và chưa có lời gọi ghi sự kiện, nên báo cáo không xem Analytics là chức năng đã triển khai. \\
Firebase AI Logic & Global Context đã được cấu hình ở vùng \code{global} trên Console. Phiên bản hiện tại chưa dùng Firebase AI SDK; Flutter gửi HTTP đến FastAPI, sau đó backend gọi Gemini bằng thư viện \code{google-genai}. \\
\bottomrule
\end{tabularx}
\end{table}

Môi trường huấn luyện được ghi nhận từ notebook Google Colab: một GPU Tesla
T4, PyTorch 2.11.0+cu130 và CUDA khả dụng. Dữ liệu được sao chép từ Google
Drive sang \code{/content} trước khi huấn luyện để giảm độ trễ đọc hàng nghìn
tệp nhỏ. Môi trường triển khai cục bộ dùng Python 3.13.7, Ultralytics 8.4.135,
PyTorch 2.13.0 và MediaPipe 1.0.1; model mặc định nằm tại
\code{backend/models/best.pt}.

Các biến cấu hình quan trọng gồm \code{MODEL\_PATH},
\code{POSE\_MODEL\_PATH}, \code{GEMINI\_API\_KEY},
\code{GEMINI\_MODEL} và \code{GEMINI\_TIMEOUT\_SECONDS}. Khóa API không được
ghi vào mã nguồn. Báo cáo chưa thực hiện benchmark trên một cấu hình cố định
gồm CPU, RAM, Android SDK và thiết bị thật, vì vậy không công bố số liệu thời
gian phản hồi theo phần cứng.

\section{Các thành phần chính}

Ứng dụng di động chịu trách nhiệm xác thực, quản lý hồ sơ cơ thể, nhận ảnh,
gửi yêu cầu phân tích, lấy gợi ý, hỏi đáp và hiển thị thử đồ theo camera.
Giao diện chính có bốn tab: Phân tích, Gợi ý, Hỏi đáp và Cá nhân. Các thao
tác mạng có trạng thái đang xử lý và lỗi; lịch sử hỏi đáp được giữ theo tài
khoản ở bộ nhớ cục bộ.

Dịch vụ API tiếp nhận JPG/PNG tối đa 10 MB, kiểm tra MIME và nội dung ảnh,
sau đó gọi checkpoint YOLO 21 lớp. Kết quả giữ nhiều hộp hợp lệ, gồm lớp gốc,
nhóm chuẩn hoá, màu chính/phụ, độ tin cậy và toạ độ. Endpoint phân tích dùng
confidence mặc định 0,25 và IoU 0,45, rồi gửi metadata cùng hồ sơ/bối cảnh
đến Gemini. Các endpoint gợi ý và hỏi đáp là luồng không ảnh, không dùng kết
quả YOLO hay dữ liệu tủ đồ.

Firebase Authentication quản lý phiên đăng nhập. Firestore lưu hồ sơ tài
khoản và danh sách hồ sơ cơ thể trong tài liệu \code{users/\{uid\}}. Security
Rules cho phép chủ tài khoản đọc/ghi subcollection
\code{wardrobe/\{itemId\}}, nhưng mã Flutter hiện chưa triển khai màn hình và
service tủ đồ; không có kho ảnh hoặc danh sách yêu thích. Dữ liệu ảnh gửi đến
backend chỉ tồn tại trong request và không được ghi ra ổ đĩa.

MediaPipe Pose Landmarker cung cấp các điểm mốc cơ thể từ ảnh hoặc video
~\cite{MediaPipePose}. Trong hệ thống, chức năng thử đồ gửi frame JPEG qua
WebSocket đến Pose Landmarker. Backend trả landmark và điểm neo; Flutter làm
mượt tư thế rồi ghép lớp trang
phục theo vùng thân trên, thân dưới, váy hoặc bàn chân. Người dùng có thể dịch
chuyển, co giãn, xoay và đặt lại lớp ảnh; frame không được lưu ở backend.

Mã huấn luyện xác nhận quy trình nạp \code{yolo26n.pt}, chuyển annotation
DeepFashion sang YOLO, kiểm tra thứ tự 21 lớp, fine-tune và lưu
\code{best.pt}. Mã triển khai nạp checkpoint này một lần khi backend khởi
động, tự chọn \code{cuda:0} nếu CUDA khả dụng và dùng CPU trong trường hợp
còn lại. Nếu Gemini thiếu khóa, hết hạn mức hoặc lỗi mạng, API vẫn trả kết
quả YOLO với \code{source=unavailable} thay vì tạo đánh giá thay thế giả.

\section{Giao diện và demo}

Luồng giao diện thực tế gồm đăng ký/đăng nhập, phân tích ảnh từ camera hoặc
thư viện, gợi ý theo bối cảnh, hỏi đáp tiếp nối, quản lý hồ sơ cá nhân và thử
trang phục theo camera. Kết quả phân tích hiển thị nhãn, màu, độ tin cậy,
điểm đánh giá và đề xuất điều chỉnh khi Gemini khả dụng.

Màn hình xác thực trong Hình~\ref{fig:ui-auth} cung cấp hai luồng tạo tài
khoản và đăng nhập bằng email/mật khẩu. Sau khi đăng nhập, người dùng truy
cập bốn nhóm chức năng chính bằng thanh điều hướng dưới màn hình.

\begin{figure}[H]
\centering
\begin{minipage}[t]{0.46\textwidth}
  \centering
  \includegraphics[width=\linewidth,height=0.57\textheight,keepaspectratio]{screenshots/ui-login.jpg}

  \small (a) Đăng nhập tài khoản
\end{minipage}\hfill
\begin{minipage}[t]{0.46\textwidth}
  \centering
  \includegraphics[width=\linewidth,height=0.57\textheight,keepaspectratio]{screenshots/ui-register.jpg}

  \small (b) Đăng ký tài khoản
\end{minipage}
\caption{Giao diện xác thực người dùng bằng email và mật khẩu.}
\label{fig:ui-auth}
\end{figure}

Khi quên mật khẩu, người dùng nhập email đã đăng ký và yêu cầu Firebase gửi
liên kết đặt lại. Hình~\ref{fig:ui-password-reset} thể hiện hộp nhập email và
thông báo xác nhận sau khi yêu cầu được gửi thành công.

\begin{figure}[H]
\centering
\begin{minipage}[t]{0.46\textwidth}
  \centering
  \includegraphics[width=\linewidth,height=0.57\textheight,keepaspectratio]{screenshots/ui-password-reset-dialog.jpg}

  \small (a) Nhập email đặt lại mật khẩu
\end{minipage}\hfill
\begin{minipage}[t]{0.46\textwidth}
  \centering
  \includegraphics[width=\linewidth,height=0.57\textheight,keepaspectratio]{screenshots/ui-password-reset-sent.jpg}

  \small (b) Xác nhận đã gửi email
\end{minipage}
\caption{Luồng yêu cầu đặt lại mật khẩu qua email.}
\label{fig:ui-password-reset}
\end{figure}

Đối với tài khoản đang đăng nhập, chức năng đổi mật khẩu yêu cầu nhập lại mật
khẩu hiện tại trước khi lưu mật khẩu mới. Hình~\ref{fig:ui-password-change}
minh hoạ biểu mẫu đổi mật khẩu và thông báo hoàn tất trên màn hình tài khoản.

\begin{figure}[H]
\centering
\begin{minipage}[t]{0.46\textwidth}
  \centering
  \includegraphics[width=\linewidth,height=0.57\textheight,keepaspectratio]{screenshots/ui-password-change-dialog.jpg}

  \small (a) Nhập mật khẩu hiện tại và mật khẩu mới
\end{minipage}\hfill
\begin{minipage}[t]{0.46\textwidth}
  \centering
  \includegraphics[width=\linewidth,height=0.57\textheight,keepaspectratio]{screenshots/ui-password-change-success.jpg}

  \small (b) Thông báo đổi mật khẩu thành công
\end{minipage}
\caption{Luồng đổi mật khẩu của tài khoản đang đăng nhập.}
\label{fig:ui-password-change}
\end{figure}

Luồng phân tích bắt đầu từ màn hình chọn ảnh hoặc chụp ảnh. Sau khi backend
xử lý, ứng dụng vẽ hộp giới hạn cùng lớp trang phục, màu và độ tin cậy như
Hình~\ref{fig:ui-detection}. Đây là bằng chứng giao diện cho thấy kết quả YOLO
đã được trả về và hiển thị trên ảnh đầu vào.

\begin{figure}[H]
\centering
\begin{minipage}[t]{0.46\textwidth}
  \centering
  \includegraphics[width=\linewidth,height=0.57\textheight,keepaspectratio]{screenshots/ui-analysis-input.jpg}

  \small (a) Chọn nguồn ảnh và hồ sơ
\end{minipage}\hfill
\begin{minipage}[t]{0.46\textwidth}
  \centering
  \includegraphics[width=\linewidth,height=0.57\textheight,keepaspectratio]{screenshots/ui-analysis-detection.jpg}

  \small (b) Hiển thị kết quả nhận diện
\end{minipage}
\caption{Luồng tiếp nhận ảnh và nhận diện trang phục bằng YOLO.}
\label{fig:ui-detection}
\end{figure}

Phần đánh giá sau nhận diện được minh hoạ ở
Hình~\ref{fig:ui-analysis-result}. Ứng dụng trình bày điểm tổng hợp, các tiêu
chí màu sắc, thời tiết, phong cách và hoàn cảnh; đồng thời nêu điểm mạnh,
điểm cần điều chỉnh và khuyến nghị phối bổ sung. Trường hợp trong ảnh cũng
cho thấy hệ thống phát hiện sự không khớp giữa nhãn chân váy và hồ sơ giới
tính Nam thay vì bỏ qua bối cảnh người dùng.

\begin{figure}[H]
\centering
\begin{minipage}[t]{0.46\textwidth}
  \centering
  \includegraphics[width=\linewidth,height=0.57\textheight,keepaspectratio]{screenshots/ui-analysis-score.jpg}

  \small (a) Điểm và tiêu chí đánh giá
\end{minipage}\hfill
\begin{minipage}[t]{0.46\textwidth}
  \centering
  \includegraphics[width=\linewidth,height=0.57\textheight,keepaspectratio]{screenshots/ui-analysis-recommendations.jpg}

  \small (b) Điểm cần điều chỉnh và khuyến nghị
\end{minipage}
\caption{Kết quả đánh giá và khuyến nghị sau khi phân tích trang phục.}
\label{fig:ui-analysis-result}
\end{figure}

Với luồng gợi ý không cần ảnh, người dùng chọn hồ sơ, nguồn thời tiết, phong
cách và hoàn cảnh. Hình~\ref{fig:ui-suggestion} cho thấy dữ liệu Open-Meteo
được hiển thị trước khi gửi yêu cầu và kết quả Gemini trả về dưới dạng mô tả
bộ đồ cùng các cách điều chỉnh phù hợp.

\begin{figure}[H]
\centering
\begin{minipage}[t]{0.46\textwidth}
  \centering
  \includegraphics[width=\linewidth,height=0.57\textheight,keepaspectratio]{screenshots/ui-suggestion-context.jpg}

  \small (a) Chọn hồ sơ và bối cảnh thời tiết
\end{minipage}\hfill
\begin{minipage}[t]{0.46\textwidth}
  \centering
  \includegraphics[width=\linewidth,height=0.57\textheight,keepaspectratio]{screenshots/ui-suggestion-result.jpg}

  \small (b) Bộ đồ và cách điều chỉnh được đề xuất
\end{minipage}
\caption{Gợi ý phối đồ dựa trên hồ sơ, thời tiết và hoàn cảnh.}
\label{fig:ui-suggestion}
\end{figure}

Hình~\ref{fig:ui-chat} minh hoạ khả năng hỏi đáp tiếp nối: câu trả lời đầu
tiên sử dụng bối cảnh mưa và nhiệt độ hiện tại, còn lượt hỏi sau tiếp tục dùng
thông tin hồ sơ và lịch sử hội thoại để tư vấn loại giày phù hợp.

\begin{figure}[H]
\centering
\begin{minipage}[t]{0.46\textwidth}
  \centering
  \includegraphics[width=\linewidth,height=0.57\textheight,keepaspectratio]{screenshots/ui-chat-initial.jpg}

  \small (a) Câu hỏi và phản hồi đầu tiên
\end{minipage}\hfill
\begin{minipage}[t]{0.46\textwidth}
  \centering
  \includegraphics[width=\linewidth,height=0.57\textheight,keepaspectratio]{screenshots/ui-chat-followup.jpg}

  \small (b) Hỏi tiếp trong cùng ngữ cảnh
\end{minipage}
\caption{Giao diện hỏi đáp phối đồ nhiều lượt với Gemini.}
\label{fig:ui-chat}
\end{figure}

Người dùng có thể tạo nhiều hồ sơ cơ thể và thay đổi hồ sơ đang hoạt động.
Các biểu mẫu trong Hình~\ref{fig:ui-profile} thu thập tên hồ sơ, chiều cao,
cân nặng và giới tính; các dữ liệu này được dùng làm bối cảnh cho phân tích và
gợi ý phối đồ.

\begin{figure}[H]
\centering
\begin{minipage}[t]{0.46\textwidth}
  \centering
  \includegraphics[width=\linewidth,height=0.57\textheight,keepaspectratio]{screenshots/ui-profile-edit.jpg}

  \small (a) Chỉnh sửa hồ sơ cơ thể
\end{minipage}\hfill
\begin{minipage}[t]{0.46\textwidth}
  \centering
  \includegraphics[width=\linewidth,height=0.57\textheight,keepaspectratio]{screenshots/ui-profile-add.jpg}

  \small (b) Tạo hồ sơ cơ thể mới
\end{minipage}
\caption{Quản lý các hồ sơ cơ thể của người dùng.}
\label{fig:ui-profile}
\end{figure}

Chức năng thử đồ cho phép người dùng căn ảnh trang phục vào khung, phóng thu,
xoay, lật và xem trước trước khi tạo lớp ghép. Hình~\ref{fig:ui-tryon-prepare}
minh hoạ bước tạo filter và kết quả ghép áo 2D theo các điểm neo tư thế. Dòng
trạng thái kết nối và FPS đồng thời cho biết frame camera đang được xử lý theo
thời gian thực.

\begin{figure}[H]
\centering
\begin{minipage}[t]{0.46\textwidth}
  \centering
  \includegraphics[width=\linewidth,height=0.57\textheight,keepaspectratio]{screenshots/ui-tryon-calibrate.png}

  \small (a) Căn ảnh trang phục và tạo filter
\end{minipage}\hfill
\begin{minipage}[t]{0.46\textwidth}
  \centering
  \includegraphics[width=\linewidth,height=0.57\textheight,keepaspectratio]{screenshots/ui-tryon-2d.png}

  \small (b) Ghép lớp áo 2D theo tư thế
\end{minipage}
\caption{Chuẩn bị ảnh trang phục và thử đồ trực tiếp ở chế độ 2D.}
\label{fig:ui-tryon-prepare}
\end{figure}

Ngoài lớp ghép 2D, ứng dụng còn cung cấp chế độ 3D bám thân. Ảnh thực nghiệm
trong Hình~\ref{fig:ui-tryon-3d} cho thấy mesh áo được biến dạng theo phần cơ
thể nhìn thấy, đồng thời người dùng vẫn có thể chuyển về 2D, xoá lớp ảnh,
chọn nhóm trang phục khác hoặc thêm ảnh đồ mới. Camera chỉ được xử lý tạm
thời và không lưu frame tại backend.

\begin{figure}[H]
\centering
\includegraphics[width=0.52\textwidth,height=0.68\textheight,keepaspectratio]{screenshots/ui-tryon-3d.png}
\caption{Kết quả thử đồ trực tiếp ở chế độ 3D bám thân.}
\label{fig:ui-tryon-3d}
\end{figure}

\clearpage

\section{Kiểm thử}

Các bộ kiểm thử tự động được chạy lại ngày 11/11/2026 và tổng hợp ở
Bảng~\ref{tab:kiemthu}. Backend dùng mock cho YOLO, Gemini, MediaPipe và API
ngoài nên không phát sinh chi phí dịch vụ; kết quả Flutter phản ánh đúng trạng
thái mã tại thời điểm biên soạn báo cáo.

\begin{table}[H]
\centering
\caption{Các nhóm ca kiểm thử của hệ thống.}
\label{tab:kiemthu}
\footnotesize
\begin{tabularx}{\textwidth}{>{\bfseries\color{daunavy}}l >{\raggedright\arraybackslash}X >{\raggedright\arraybackslash}X >{\color{daumaroon}}l}
\toprule
\rowcolor{gray100}\bfseries\color{daunavy}Nhóm & \bfseries\color{daumaroon}Đầu vào/thao tác & \bfseries\color{daunavy}Kết quả mong đợi & \bfseries\color{daumaroon}Thực tế \\
\midrule
Backend & validation ảnh, YOLO, Gemini, thời tiết, pose và WebSocket & route trả dữ liệu/lỗi đúng schema, không lưu frame & 104/104 đạt \\
Flutter & model, service, widget và điều hướng ứng dụng & toàn bộ test tự động đạt & 139 đạt, 6 lỗi \\
Lỗi giao diện & hồ sơ cá nhân, gợi ý không ảnh, đổi giới tính, thời tiết và chuyển tab & nội dung và trạng thái khớp kỳ vọng test & chưa đạt \\
Thử đồ & ánh xạ camera, làm mượt pose, filter, mesh và thao tác lớp ảnh & tính toán ổn định trên test tự động & đạt trong bộ test \\
\bottomrule
\end{tabularx}
\end{table}

Sáu lỗi Flutter tập trung ở các test widget đã lệch với giao diện hiện tại:
hiển thị tên tại màn Cá nhân; trạng thái giới tính khi yêu cầu gợi ý đang chờ;
ba trường hợp nội dung Gợi ý/Open-Meteo; và luồng mở Gợi ý từ thanh tab. Đây
là lỗi cần đồng bộ giữa giao diện và test trước khi tuyên bố toàn bộ ứng dụng
đạt kiểm thử hồi quy.

